% Figure: coupled-tank salt masses over time, showing the fast
% antisymmetric mode and the slow symmetric mode.
\begin{figure}[htbp]
\centering
\begin{tikzpicture}
\begin{axis}[
  width=11cm, height=6cm,
  xmin=0, xmax=5, ymin=0, ymax=10.5,
  xlabel={$t$}, ylabel={salt mass},
  xtick={0,1,2,3,4,5},
  ytick={0,2,4,6,8,10},
  label style={font=\small}, tick label style={font=\scriptsize},
  grid=major, grid style={enggray!25},
  legend style={font=\scriptsize, at={(0.97,0.97)}, anchor=north east, draw=enggray!60},
]
% x1 = 5 e^{-t} + 5 e^{-3t}
\addplot[engblue, very thick, domain=0:5, samples=140] {5*exp(-x) + 5*exp(-3*x)};
\addlegendentry{$x_{1}(t)=5e^{-t}+5e^{-3t}$}
% x2 = 5 e^{-t} - 5 e^{-3t}
\addplot[engred, very thick, domain=0:5, samples=140] {5*exp(-x) - 5*exp(-3*x)};
\addlegendentry{$x_{2}(t)=5e^{-t}-5e^{-3t}$}
% slow symmetric envelope 5 e^{-t}
\addplot[enggray, densely dashed, domain=0:5, samples=100] {5*exp(-x)};
\addlegendentry{slow mode $5e^{-t}$}
% mark x2 peak at t = (ln3)/2 approx 0.549, x2 = 5(e^-0.549 - e^-1.648) = 5(0.577-0.192)=1.92
\draw[densely dashed, enggray] (axis cs:0.549,0) -- (axis cs:0.549,1.92);
\node[font=\scriptsize, engred, anchor=south west] at (axis cs:0.58,1.95) {$x_{2}$ peak};
\end{axis}
\end{tikzpicture}
\caption[Coupled-tank modal decay]{Salt masses in the two recirculating
tanks, starting with all salt in tank~1. Tank~2 rises from zero as the
fast antisymmetric mode ($e^{-3t}$) floods salt across the imbalance,
peaks at $t = \tfrac12\ln 3 \approx 0.55$, and then both tanks empty
together along the slow symmetric mode $e^{-t}$ (dashed) once the
imbalance has relaxed. The eigen-decomposition separates the dynamics
into two independent modes with distinct decay rates, and the engineer
reads the two timescales off the eigenvalues $-1$ and $-3$ rather than
off the coupled equations.}
\label{fig:vol02:ch07:coupled-tanks}
\end{figure}
